\chapter{Epílogo y Glosario Amigable}
\label{cap:epilogo_glosario}

\section{Epílogo: La matemática viva que protege al mundo}

¡Felicidades por llegar hasta aquí!

Al iniciar este libro, tal vez pensabas que la criptografía post-cuántica era un territorio impenetrable, reservado únicamente para matemáticos con décadas de experiencia o investigadores de laboratorios nucleares.

Sin embargo, a lo largo de estas páginas has descubierto una verdad inspiradora:
\begin{enumerate}
  \item Los conceptos fundamentales no son monstruos incomprensibles; son \textbf{ideas geométricas puras y elegantes}.
  \item Una cuadrícula de puntos regulares en el plano que aprendiste a graficar en tu cuaderno de colegio es el germen de un \textbf{retículo multidimensional}.
  \item La aritmética de las horas de un reloj que utilizas para saber cuándo empieza tu clase es la misma \textbf{aritmética modular $\mathbb{Z}_q$} que confina los secretos digitales.
  \item El juego de mezclar pintura o susurrar en un concierto ruidoso es exactamente lo que hace el problema \textbf{LWE}, donde un error deliberado frustra los métodos convencionales de resolución.
  \item Los polinomios que sumas en el pizarrón permiten que las matrices viajen livianas por las redes de fibra óptica en \textbf{Module-LWE}.
\end{enumerate}

Las dos grandes primitivas que acabas de comprender ---\textbf{ML-KEM} para el candado de confidencialidad y \textbf{ML-DSA} para el sello de autenticidad--- no son teorías de ciencia ficción archivadas en una biblioteca: son los estándares oficiales del Instituto Nacional de Estándares y Tecnología (NIST) aprobados en agosto de 2024 para blindar toda la infraestructura informática del planeta durante los próximos cincuenta años.

Y son exactamente las primitivas matemáticas que dan vida al artefacto de investigación desarrollado en la Universidad Mayor de San Andrés (UMSA) para garantizar que los votos ciudadanos sean anónimos, incorruptibles y eternamente seguros frente a cualquier máquina cuántica del futuro.

\section{Glosario Amigable para Primer Semestre}

Guarda este glosario como una brújula rápida para tus estudios:

\begin{description}
  \item[\textbf{Aritmética Modular ($\mathbb{Z}_q$):}] Sistema numérico circular finito donde los cálculos dan la vuelta al llegar a un valor tope $q$, quedándose con el residuo de la división entera.
  
  \item[\textbf{Base de un Retículo:}] Conjunto mínimo de vectores cuyas combinaciones lineales enteras generan todos y cada uno de los puntos de la red. Una base ``buena'' tiene vectores cortos y casi ortogonales; una base ``mala'' tiene vectores largos y distorsionados.
  
  \item[\textbf{CVP (\textit{Closest Vector Problem}):}] Problema geométrico que consiste en encontrar el punto del retículo más cercano a un punto que fue perturbado por un pequeño ruido en el espacio.
  
  \item[\textbf{Firma Digital:}] Equivalente criptográfico moderno de un sello de cera real: garantiza que el mensaje no fue modificado (integridad) y que proviene legítimamente del autor (autenticidad).
  
  \item[\textbf{KEM (\textit{Key Encapsulation Mechanism}):}] Protocolo criptográfico diseñado para acordar una clave secreta simétrica compacta entre dos personas a través de un canal público no confiable.
  
  \item[\textbf{LWE (\textit{Learning With Errors}):}] Problema computacional introducido por Oded Regev en 2005. Consiste en resolver un sistema lineal ruidoso $\mathbf{b} = \mathbf{A}\mathbf{s} + \mathbf{e} \pmod q$. Es inmune a computadoras clásicas y cuánticas.
  
  \item[\textbf{ML-DSA (FIPS 204):}] Algoritmo de firma digital post-cuántico basado en retículos modulares (Module-LWE y Module-SIS). Anteriormente conocido como *CRYSTALS-Dilithium*.
  
  \item[\textbf{ML-KEM (FIPS 203):}] Mecanismo de encapsulamiento de claves post-cuántico basado en Module-LWE. Anteriormente conocido como *CRYSTALS-Kyber*.
  
  \item[\textbf{Module-LWE (M-LWE):}] Variante algebraica de LWE donde las coordenadas de los vectores y las entradas de las matrices son polinomios de un anillo $R_q = \mathbb{Z}_q[x]/(x^n+1)$, ofreciendo tamaño compacto y cálculo ultrasónico.
  
  \item[\textbf{Muestreo con Rechazo (\textit{Rejection Sampling}):}] Técnica estadística empleada en firmas digitales de retículos donde los intentos que podrían filtrar información sobre la clave secreta se descartan y recalculan antes de publicarse.
  
  \item[\textbf{Retículo (\textit{Lattice}):}] Conjunto discreto e infinito de puntos en un espacio vectorial generado por todas las combinaciones lineales enteras de un conjunto de vectores base linealmente independientes.
  
  \item[\textbf{SIS (\textit{Short Integer Solution}):}] Problema de hallar un vector no nulo de números enteros muy pequeños $\mathbf{z}$ tal que $\mathbf{A}\mathbf{z} \equiv \mathbf{0} \pmod q$.
\end{description}

\vspace{0.8cm}
\begin{tcolorbox}[colback=BoxBgBlue, colframe=PQCMainBlue, arc=2mm, title={\sffamily\bfseries Para Seguir Aprendiendo}]
\sffamily\small
Si quieres profundizar en este fascinante viaje:
\begin{itemize}
  \item Revisa las especificaciones oficiales del NIST: \textbf{FIPS 203} (ML-KEM) y \textbf{FIPS 204} (ML-DSA) publicadas en 2024.
  \item Consulta los artículos fundacionales de Oded Regev (LWE, 2005) y Vadim Lyubashevsky (Firmas con rechazo, 2009).
  \item Explora el código fuente y benchmarks del proyecto de grado en el repositorio \texttt{src/pqc} de la Carrera de Informática de la UMSA.
\end{itemize}
\end{tcolorbox}
